\section{Asynchronous Checking}
\label{sec:async}

Aptos executes blocks of transactions in parallel using \emph{Block-STM}~\cite{BlockSTM}, an optimistic concurrency-control protocol in which a pool of worker threads speculatively runs transactions out of order and re-validates them against later-published writes; incarnations that lose a race are aborted and replayed.
In practice most blocks do not saturate the worker pool, and aborted-and-restarted incarnations leave threads idle for non-trivial windows, so a substantial slice of worker capacity remains unused during normal operation.

We use that idle capacity to run the runtime safety checks \emph{asynchronously}.
During execution the interpreter trails a compact \emph{execution trace} for each transaction, containing no values, only control-flow shape.
%: a count of successfully executed instructions, a fingerprint hash of the executed opcodes (used as an integrity tag), the bit-vector of conditional-branch outcomes, and the list of dynamically dispatched calls (closures and entry points; static calls are reconstructable from the bytecode and are not recorded).
%The trace contains no values, only control-flow shape, so it stays small.

When a transaction's commit-time incarnation finalizes, an idle worker replays its trace by abstractly interpreting the bytecode along the recorded path --- following the recorded branches, dispatching the recorded closures --- and runs the type and ability checks of \S\ref{sec:type} and \S\ref{sec:ability} on the fully-substituted types.
% Because the trace records only \emph{successful} instructions, gas need not be re-charged; the async checker uses an unmetered gas meter.
%, and the recorded fingerprint catches any divergence between the original and the replayed instruction streams.
A violation found during replay aborts the block before its writes become observable to subsequent blocks.
Because asynchronous checks are run on successful transactions only, speculative executions and aborted transactions do not pay the cost of these checks.
Because the checks are almost always expected to pass on the static verifier admitted bytecode, except in the low probability event of a latent bug identified in the verifier being exploited, computation done while executing a transaction is not wasted because of failed checks.

%%% Local Variables:
%%% mode: latex
%%% TeX-master: "main"
%%% End:
